\subsection{Vision Transformer Experiments}
\label{app:vit}

\paragraph{Relation to the original vision experiments.}
Uni-LoRA \citep{li2025unilora} already evaluates ViT-Base and ViT-Large on eight image-classification datasets, with five runs per configuration (its Section~4.4 and Table~5). We extend this line of evaluation to ProLoSA using the completed ViT-B/16 runs. Our data subsets, optimization settings, and seed coverage differ from that study, so all LoRA and Uni-LoRA scores in Table~\ref{tab:vit} come from the same local experiment suite as ProLoSA. They are not copied from the original paper. No completed ProLoSA ViT-Large results are available in this suite.

\paragraph{Model and data.}
We use the ImageNet-21k-pretrained ViT-B/16 checkpoint (\path{google/vit-base-patch16-224-in21k}) and apply rank-4 adapters to query and value projections in all 12 Transformer blocks. All other backbone parameters are frozen, and a randomly initialized classification head is trained for every method. The three limited-data settings use 1,000 training images from CIFAR-100, Food-101, and DTD, respectively. Sampling is stratified by class with subset seed 42, giving 10 images per class for CIFAR-100 and approximately balanced allocations for Food-101 and DTD. CIFAR-100 additionally uses 5,000 images and the full 50,000-image training split. These are custom stratified subsets, not the standard VTAB-1k protocol. Evaluation uses the complete CIFAR-100 test split (10,000 images), Food-101 validation split (25,250), and DTD test split (1,880), respectively. Training applies random resized cropping to $224\times224$, horizontal flipping, and pretrained-model normalization; evaluation uses resizing, center cropping, and the same normalization.

\paragraph{Budgets and parameter accounting.}
The rank-4 LoRA factor space contains $D=12\times2\times(768\times4+4\times768)=147{,}456$ parameters. Uni-LoRA uses $d=24{,}600$, while ProLoSA divides the same budget into $d+K=24{,}600$ (Table~\ref{tab:vit_allocations}). The shared head adds $769C$ parameters for $C$ classes: 76,900 for CIFAR-100, 77,669 for Food-101, and 36,143 for DTD. Thus the total active parameter counts for either compressed method are 101,500, 102,269, and 60,743, respectively; LoRA totals are 224,356, 225,125, and 183,599. Adapter-only compression is approximately $6\times$; this factor does not apply to the totals including the head. ProLoSA counts the $K$ selected sparse coefficients even though they are frozen during warmup. The raw result files count parameters before activation and therefore contain only $d$ in their ProLoSA adapter field; we recover $d+K$ from the logged configuration and verify $K$ against the activation log. These counts describe active degrees of freedom, not allocated tensor storage or optimizer memory.

\begin{table}[t]
\centering\small
\begin{tabular}{lrr@{\qquad}lrr}
\toprule
\multicolumn{3}{c}{Main budget: $d+K=24{,}600$} & \multicolumn{3}{c}{Supplementary budget: $d+K=72{,}000$} \\
\cmidrule(lr){1-3}\cmidrule(lr){4-6}
Method & $d$ & $K$ & Method & $d$ & $K$ \\
\midrule
Uni-LoRA & 24,600 & 0 & Uni-LoRA & 72,000 & 0 \\
ProLoSA ($2:1$) & 16,400 & 8,200 & ProLoSA ($2:1$) & 48,000 & 24,000 \\
ProLoSA ($4:1$) & 19,680 & 4,920 & ProLoSA ($4:1$) & 57,600 & 14,400 \\
ProLoSA ($8:1$) & 21,867 & 2,733 & ProLoSA ($8:1$) & 64,000 & 8,000 \\
ProLoSA ($12:1$) & 22,708 & 1,892 & ProLoSA ($12:1$) & 66,462 & 5,538 \\
\bottomrule
\end{tabular}
\caption{Exact compressed-adapter budgets for ViT-B/16. Ratios are approximate when integer rounding is required. The task-specific classification head is additional and shared across methods.}
\label{tab:vit_allocations}
\end{table}

\paragraph{Optimization and support selection.}
All runs use AdamW with $(\beta_1,\beta_2)=(0.9,0.999)$, $\epsilon=10^{-8}$, zero weight decay, BF16, batch size 64, and no gradient accumulation. The 1,000-image settings train for 20 epochs (320 optimizer steps); CIFAR-100 with 5,000 and 50,000 images trains for 10 epochs (790 and 7,820 steps). LoRA uses $\alpha=r=4$ and adapter learning rate $10^{-3}$; Uni-LoRA and ProLoSA use a latent-vector learning rate of $4\times10^{-3}$. The head learning rate is $3\times10^{-3}$ for every method. Adapter dropout is zero. The latent/head learning rates follow linear decay with 10\% scheduler warmup. ProLoSA uses latent initialization bound 0.02, a score-free warmup of 32 steps for the 1,000-image settings or 128 steps otherwise, and a one-step SNIP scoring window. Its support is activated after step 33 or 129, respectively, with zero-initialized sparse coefficients and cleared optimizer state. The sparse learning rate is $0.2$ times the latent rate ($8\times10^{-4}$), with linear decay after activation. This support-selection schedule differs from the 128-step scoring window used in the earlier language experiments.

\paragraph{Reporting and limitations.}
We report \emph{final-model} top-1 accuracy, using the \texttt{final\_accuracy} field rather than \texttt{best\_accuracy}, which is the maximum over repeated evaluations on the same held-out split. The full method-by-setting comparison is available only for training seed 42, and Table~\ref{tab:vit} retains all four ProLoSA ratios. Only LoRA on CIFAR-100 with 1,000 images has additional completed seeds: its final accuracies for seeds 42/43/44 are 86.91/86.57/86.44. We therefore do not attach multi-seed error bars to the compressed-method comparison. No independent validation split was used to choose between the available budget/learning-rate recipes, so the highlighted recipe is an exploratory result. Full-data CIFAR-100 improves by 0.23 points over Uni-LoRA at $d:K=4:1$, but the four ProLoSA ratios lose 0.09--0.25 points at 1,000 images and 0.19--0.42 points at 5,000 images. The small Food-101 and DTD differences and the lack of repeated ProLoSA seeds preclude a consistent-advantage claim. The changing epoch count and number of updates across data sizes also prevent a controlled test of the sample-size law.

\paragraph{Supplementary 72,000-parameter recipe.}
Table~\ref{tab:vit_72k} retains the second complete single-seed sweep. It uses $d=72{,}000$ for Uni-LoRA and $d+K=72{,}000$ for ProLoSA, with latent learning rate $10^{-2}$, sparse learning rate $2\times10^{-3}$, and head learning rate $10^{-2}$ on CIFAR-100 or $5\times10^{-2}$ on Food-101 and DTD. LoRA's adapter learning rate remains $10^{-3}$; data, epochs, and support-selection timing are unchanged. Since the learning rates change together with the compressed budget (including the head rate for LoRA), differences from Table~\ref{tab:vit} do not isolate a budget effect. At this recipe, ProLoSA's $2:1$ allocation exceeds Uni-LoRA on the four limited-data settings but loses on full-data CIFAR-100; LoRA remains stronger than this allocation on Food-101, DTD, and the larger CIFAR-100 subsets. This recipe dependence reinforces the need for independent tuning and repeated seeds.

\begin{table}[t]
\centering\small
\setlength{\tabcolsep}{5pt}
% Generated from final_accuracy, seed 42; do not edit.
\begin{tabular}{lrrrrr}
\toprule
 & \multicolumn{3}{c}{1,000 training images} & \multicolumn{2}{c}{CIFAR-100} \\
\cmidrule(lr){2-4}\cmidrule(lr){5-6}
Method & CIFAR-100 & Food-101 & DTD & 5,000 & Full \\
\midrule
LoRA ($r=4$) & 85.22 & 73.52 & 68.83 & 89.32 & 91.78 \\
Uni-LoRA & 85.10 & 72.68 & 66.70 & 88.24 & 91.51 \\
ProLoSA ($2:1$) & 85.38 & 72.77 & 67.93 & 88.96 & 91.43 \\
ProLoSA ($4:1$) & 85.08 & 71.95 & 68.19 & 88.55 & 91.80 \\
ProLoSA ($8:1$) & 84.79 & 72.48 & 68.09 & 88.60 & 91.57 \\
ProLoSA ($12:1$) & 84.98 & 72.37 & 67.93 & 88.72 & 91.76 \\
\bottomrule
\end{tabular}

\caption{Supplementary ViT-B/16 recipe: final-model accuracy (\%, seed 42), with 72,000 active adapter parameters for both compressed methods and 147,456 for LoRA, plus the shared classification head. All available ProLoSA ratios are retained. Learning rates differ from the main vision table as specified in the text.}
\label{tab:vit_72k}
\end{table}
